% id: 115-46
% book: 쎈B 공통수학2 · 08 함수
% section: 유형 11 $f\circ g=g\circ f$인 경우
% figure: none
% answer: 
% answer_source: 
% variant_level: 0 (원본 전사)
% 이 파일은 build.mjs 가 items.json 에서 만든다. 고칠 때는 items.json 을 고친다.
\smash{\raisebox{1.5mm}{\dmkichul{쎈B 공통수학2 115쪽 46번 · 서술형}}}
두 함수 $f(x)=-2x+6$, $g(x)=ax+b$가 $f\circ g=g\circ f$를 만족시킬 때, 함수 $y=g(x)$의 그래프는 $a$의 값에 관계없이 항상 한 점을 지난다. 이 점의 좌표를 구하시오.\dmcondfill{$a$, $b$는 상수이다.}{}
